\chapter{Chart proofs of rectifiable area and coarea}\label{sec:rectifiable-chart-integration}
This section develops the chart mechanism behind \cite{Sim18}, Chapter 3, equations (2.4)--(2.9) and Remark 2.10, printed pp.~92--94. The Euclidean area and coarea formulas remain external inputs. The tangential differential and its independence, the $C^1$ decomposition, and the tangent theorem are supplied by the preceding development. Throughout, $U\subset\mathbb R^n$ and the coordinate domains are open, and $1\leq m\leq n$. Unless otherwise specified, all weights are nonnegative and all integrals may be infinite.

\begin{lemma}[Area Jacobians under a change of tangent coordinates]\label{lem:tangent-area-jacobian-chain}
Let $A:\mathbb R^m\to\mathbb R^n$ be injective and linear, let $T=A(\mathbb R^m)$, and let $L:T\to\mathbb R^q$ be linear with $q\geq m$. Then
\[
 J_m(LA)=J_mL\,J_mA.
\]
Here $J_mL$ uses the Euclidean metric restricted to $T$; either side can vanish if $L$ loses rank.
\end{lemma}
\begin{proof}
Choose an isometry $Q:\mathbb R^m\to T$, and write $A=QR$ with an invertible $m\times m$ matrix $R$. The Gram matrix of $LA$ is $R^*(LQ)^*(LQ)R$, whose determinant is $\det(R)^2\det((LQ)^*(LQ))$. Also $J_mA=|\det R|$ and $J_mL=\sqrt{\det((LQ)^*(LQ))}$. Taking square roots proves the identity, including a zero Gram determinant. Orthogonal choice of $Q$ does not change any of these quantities.
\end{proof}

\begin{theorem}[Weighted area on a rectifiable domain]\label{thm:rectifiable-weighted-area-local}
Let $M\subset U\subset\mathbb R^n$ be measurable, countably $m$-rectifiable and locally finite in $\mathcal H^m$ measure, let $f:U\to\mathbb R^q$ be locally Lipschitz with $q\geq m$, and let $h:M\to[0,\infty]$ be measurable. Then
\[
 \int_M h(x)J_m(d^Mf_x)\,d\mathcal H^m(x)
 =\int_{\mathbb R^q}\sum_{x\in M\cap f^{-1}(y)}h(x)\,d\mathcal H^m(y).
\]
The fiber sum has a measurable representative. This is \cite{Sim18}, equations (2.4)--(2.6), derived here from the Euclidean weighted area formula. Products with a zero Jacobian use the convention $0\cdot\infty=0$.
\end{theorem}
\begin{proof}
Choose a Borel representative for $M$ and a Borel representative for $h$ relative to its sigma-finite restricted Hausdorff measure. This changes them only on an $\mathcal H^m$-null subset; the image of that subset is null by the Lipschitz bound. Partition this representative of $M$, apart from an $\mathcal H^m$-null set, into disjoint Borel pieces $M_i$ lying in injective regular $C^1$ charts $F_i:V_i\to U$ as in Lemma \ref{lem:finite-area-c1-patches}. Subdivide so the chart and the relevant restrictions of $f$ are Lipschitz on each piece's neighborhood. Put $A_i=F_i^{-1}(M_i)$ and $g_i=f\circ F_i$. At almost every point of $A_i$, Theorem \ref{thm:rectifiable-tangential-differential} gives $Dg_i=d^Mf\circ DF_i$. Lemma \ref{lem:tangent-area-jacobian-chain} therefore gives
\[
 J_mDg_i=J_m(d^Mf)J_mDF_i.
\]
Apply Proposition \ref{prop:weighted-area} to $g_i$ with weight $\chi_{A_i}(h\circ F_i)$. Apply the same formula to the injective chart $F_i$ to change the left integral to integration on $M_i$. The resulting right sum is over precisely the points of $M_i\cap f^{-1}(y)$, because $F_i$ is injective. Sum over $i$ using monotone convergence; the pieces are disjoint, so each point is counted once. This construction also gives measurability of the resulting sum modulo null sets.

The omitted Hausdorff-null set has Hausdorff-null image under $f$, by the elementary Lipschitz bound on a countable local cover. Its fiber contribution therefore vanishes for almost every target point. On the left it contributes zero as well. This completes the formula on $M$ and permits infinite values by monotone approximation of $h$.
\end{proof}

\begin{lemma}[A Hausdorff-null set has null lower-dimensional fibers]\label{lem:null-set-lipschitz-fibers}
Let $1\leq k<m$, let $f:U\subset\mathbb R^n\to\mathbb R^k$ be locally Lipschitz, and let $Z\subset U$ have $\mathcal H^m(Z)=0$. Then
\[
 \mathcal H^{m-k}(Z\cap f^{-1}(y))=0
 \quad\text{for }\mathcal L^k\text{-almost every }y.
\]
No rectifiability of $Z$ is assumed. This is the null-set part of the Lipschitz slicing estimate in \cite{Sim18}, Chapter 2, Theorem 1.10, pp.~49--51, proved directly by diameter covers.
\end{lemma}
\begin{proof}
Localize to $Z$ in a compact subset of a ball where $f$ has Lipschitz constant $L$ and whose closure lies in $U$. A countable such localization covers the original set. Put $d=m-k>0$. For each integer $l$, Hausdorff-nullness gives a cover by closed balls $Q_{li}$ of positive radii $r_{li}<2^{-l}$, all in the chosen neighborhood, with $\sum_i r_{li}^m<2^{-l}$. One obtains such balls from diameter covers by choosing centers in $Z$ and enlarging radii by a fixed factor; countably many zero-diameter cover members may be discarded, since their intersections with every fiber have zero $\mathcal H^d$ measure.

The sets $f(Q_{li})$ are compact, so
\[
 F_l(y)=\sum_i\omega_d r_{li}^d\chi_{f(Q_{li})}(y)
\]
is measurable. Each image is contained in a $k$-ball of radius $Lr_{li}$ about the image of its center. Hence
\[
 \int F_l\,dy\leq\omega_d\omega_kL^k\sum_i r_{li}^{d+k}
 <\omega_d\omega_kL^k2^{-l}.
\]
Tonelli implies $\sum_lF_l(y)<\infty$ for almost every $y$, so $F_l(y)\to0$ there. The corresponding selected balls cover $Z\cap f^{-1}(y)$ and have diameters less than $2^{1-l}$. Their normalized Hausdorff cost is at most $F_l(y)$. Thus $\mathcal H^d_{2^{1-l}}(Z\cap f^{-1}(y))\leq F_l(y)$, and sending $l\to\infty$ gives zero outer Hausdorff measure. The countable exceptional sets from localization and the discarded countable cover members cause no additional problem. This proof also avoids presupposing measurability of an arbitrary fiber measure function.
\end{proof}

\begin{theorem}[Regular geometry of almost every Euclidean Lipschitz fiber]\label{thm:euclidean-lipschitz-fiber-geometry}
Let $V\subset\mathbb R^m$ be open and $g:V\to\mathbb R^k$ be locally Lipschitz, with $1\leq k<m$. For almost every $y\in\mathbb R^k$, the fiber $N_y=g^{-1}(y)$ is countably $(m-k)$-rectifiable and locally finite in $\mathcal H^{m-k}$ measure in $V$. At $\mathcal H^{m-k}$-almost every point of each such fiber, $g$ is differentiable, has derivative of rank $k$, and
\[
 T_xN_y=\ker Dg(x).
\]
The proof uses the cited Euclidean coarea formula and the local $C^1$ approximation theorem.
\end{theorem}
\begin{proof}
On a countable family of relatively compact balls covering $V$, extend the restriction of each component of $g$ to a global Lipschitz function by Proposition \ref{prop:mcshane-extension}. Corollary \ref{cor:vector-c1-lusin-approximation}, with exceptional measures tending to zero, supplies countably many $C^1$ maps $G_i$ which agree with $g$ and $Dg$ on sets whose union covers each ball up to a Lebesgue-null set. Collect those null sets and the nondifferentiability set into $Z$. The Euclidean coarea formula, Theorem \ref{thm:coarea-formula}, implies that $Z$ meets almost every fiber in an $\mathcal H^{m-k}$-null set. The same formula applied to the set where $\operatorname{rank}Dg<k$ gives a null contribution to almost every fiber, since its $k$-Jacobian is zero.

At every remaining point where $g=G_i$ and $Dg=DG_i$, the derivative has rank $k$. The implicit function theorem makes the regular level $G_i^{-1}(y)$ a $C^1$ submanifold of dimension $m-k$ near that point, with tangent $\ker DG_i$. These regular level manifolds admit countably many charts. For a fixed good $y$, they cover the actual fiber up to a null set, proving countable rectifiability. The Euclidean coarea formula on a countable compact exhaustion gives local finite fiber measure for almost every $y$, because $J_kg$ is bounded on each compact neighborhood and that neighborhood has finite volume.

At almost every point of the common part of the fiber and each regular level patch, the tangent agreement in Theorem \ref{thm:rectifiable-forward-tangents} identifies the fiber tangent with the patch tangent. On the agreement subset under consideration $DG_i=Dg$, so that tangent is $\ker Dg$. Remove the countable union of all exceptional subsets. This proves the simultaneous geometry assertion. The coarea input is a scalar integral identity here; the fiber rectifiability and tangent identification have been derived from it and the $C^1$ approximation.
\end{proof}

\begin{lemma}[The Jacobian identity separating a fiber from its normal directions]\label{lem:coarea-chart-jacobian-identity}
Let $A:\mathbb R^m\to\mathbb R^n$ be injective and linear, and let $B:\mathbb R^m\to\mathbb R^k$ be linear of rank $k$, where $1\leq k<m$. Set $T=A(\mathbb R^m)$ and define $L:T\to\mathbb R^k$ by $LA=B$. With $K=\ker B$,
\[
 J_mA\,J_kL=J_kB\,J_{m-k}(A|_K).
\]
The $k$-Jacobians use the coarea convention $J_kL=\sqrt{\det(LL^*)}$. The $m$- and $(m-k)$-Jacobians use the area convention. This is the coordinate identity behind \cite{Sim18}, equation (2.8).
\end{lemma}
\begin{proof}
Choose orthonormal coordinates on $\mathbb R^m=K\oplus K^\perp$. Then $B=(0,C)$ with invertible $k\times k$ matrix $C$. Write the positive definite metric $G=A^*A$ in corresponding blocks $G_{11},G_{12},G_{21},G_{22}$. The tangential coarea matrix is $BG^{-1}B^*=C(G^{-1})_{22}C^*$, because $A$ identifies $T$ with this metric on the domain. The block inverse formula and the Schur determinant formula give
\[
 \det G\,\det((G^{-1})_{22})=\det G_{11}.
\]
Indeed both formulas use the positive definite Schur complement $G_{22}-G_{21}G_{11}^{-1}G_{12}$. Hence
\[
 (J_mA\,J_kL)^2
 =\det G\,\det(BG^{-1}B^*)
 =\det(C)^2\det G_{11}.
\]
The factors on the last line are $(J_kB)^2$ and $(J_{m-k}(A|_K))^2$. Taking positive square roots proves the identity. The proof makes explicit which metric enters the tangential coarea Jacobian; replacing it by $BB^*$ before accounting for the chart metric would be incorrect.
\end{proof}

\begin{theorem}[Rectifiable coarea and the tangent of almost every fiber]\label{thm:rectifiable-coarea-chart-proof}
Let $M\subset U\subset\mathbb R^n$ be measurable, countably $m$-rectifiable and locally finite in $\mathcal H^m$ measure. Let $f:U\to\mathbb R^k$ be locally Lipschitz, $1\leq k<m$, and let $h:M\to[0,\infty]$ be measurable. Then
\[
 \int_M hJ_k(d^Mf)\,d\mathcal H^m
 =\int_{\mathbb R^k}\int_{M\cap f^{-1}(y)}h\,d\mathcal H^{m-k}\,dy.
\]
For almost every $y$, the fiber is countably $(m-k)$-rectifiable and locally finite in $\mathcal H^{m-k}$ measure. At almost every point of such a fiber, $d^Mf$ has rank $k$ and
\[
 T_x(M\cap f^{-1}(y))=\ker d^Mf_x\subset T_xM.
\]
This is the content of \cite{Sim18}, Chapter 3, equation (2.8), its weighted extension (2.9), and Remark 2.10(2), pp.~93--94, with the fiber tangent step supplied explicitly.
\end{theorem}
\begin{proof}
First choose Borel representatives of $M$ and $h$ modulo an $\mathcal H^m$-null set; Lemma \ref{lem:null-set-lipschitz-fibers} shows that this changes almost every fiber only by an $\mathcal H^{m-k}$-null set. Choose disjoint Borel chart pieces $M_i$ with injective $C^1$ charts $F_i:V_i\subset\mathbb R^m\to U$, and put $A_i=F_i^{-1}(M_i)$, $g_i=f\circ F_i$. Work first in one chart, subdividing into relatively compact base balls where the relevant maps are Lipschitz. The Borel choice makes every coordinate piece and its intersections with the closed fibers measurable for the relevant Hausdorff measures. Theorem \ref{thm:euclidean-lipschitz-fiber-geometry} gives a rectifiable, locally finite fiber of $g_i$ at almost every level, with tangent $\ker Dg_i$ almost everywhere. Its measurable subset $A_i\cap g_i^{-1}(y)$ has the same tangent almost everywhere on that subset by Corollary \ref{cor:rectifiable-common-tangents}.

Apply Theorem \ref{thm:rectifiable-weighted-area-local} to the injective map $F_i$ on this base fiber, with weight $h\circ F_i$. It transforms the fiber integral into
\[
 \int_{A_i\cap g_i^{-1}(y)}
 (h\circ F_i)\,J_{m-k}(DF_i|_{\ker Dg_i})\,d\mathcal H^{m-k}.
\]
The kernel expression is set to zero on the rank-deficient or nondifferentiable set, which has null fiber measure for almost every level. The ambient-null set where $Dg_i=d^Mf\circ DF_i$ or the prescribed tangential representative fails is also removed: Euclidean coarea makes its intersection with almost every base fiber null. It is measurable on the full-rank set: the orthogonal kernel projection is $I-(Dg_i)^*(Dg_i(Dg_i)^*)^{-1}Dg_i$, a continuous matrix expression in a full-rank derivative.

Integrate the last display in $y$ and apply the Euclidean coarea formula, extended to nonnegative weights by simple functions and monotone convergence. Lemma \ref{lem:coarea-chart-jacobian-identity}, with $A=DF_i$, $B=Dg_i$ and $L=d^Mf$, identifies the resulting integrand with $(h\circ F_i)J_mDF_iJ_k(d^Mf)$. At points of rank less than $k$, both coarea Jacobians vanish, so this equality of integrands still holds under the zero convention. The chart area formula then gives $\int_{M_i}hJ_k(d^Mf)\,d\mathcal H^m$. Summing over the disjoint pieces proves the integral identity on their union.

The omitted $\mathcal H^m$-null subset of $M$ meets almost every fiber in an $\mathcal H^{m-k}$-null set by Lemma \ref{lem:null-set-lipschitz-fibers}. It affects neither side. The fiber is therefore covered modulo a null set by countably many Lipschitz images of the rectifiable base fibers, proving its rectifiability. Apply the proved formula with indicators of a countable compact exhaustion. Local Lipschitz bounds give $J_k(d^Mf)\leq(\operatorname{Lip}f)^k$ on each neighborhood and the unweighted area of $M$ is locally finite. Thus the fiber measures are locally finite for almost every level.

Finally each regular base fiber is covered modulo null sets by the $C^1$ regular level pieces used in Theorem \ref{thm:euclidean-lipschitz-fiber-geometry}. Their images under the injective regular chart are $C^1$ submanifolds with tangent $DF_i(\ker Dg_i)=\ker d^Mf$. The tangent transfer theorem identifies the actual fiber's tangent with these piece tangents at almost every point. The countable union of exceptional subsets is null. This proves the rank and tangent assertions simultaneously. The fundamental Euclidean area and coarea formulas remain external; no circular appeal to the rectifiable coarea theorem has entered this proof.
\end{proof}
